122
11 Multiple Integrals
and the discontinuities of $\chi_E$ form the set $\partial E$, we find by Lebesgue's criterion that
the measure just introduced is defined only for admissible sets.
Thus admissible sets, and only admissible sets, are measurable in the sense of
Definition 3.
Let us now ascertain the geometric meaning of $\mu(E)$. If $E$ is an admissible set
then
$$\mu(E) = \int_I \chi_E(x) dx = \int_{\underline{I}} \chi_E(x)dx = \int_{\overline{I}} \chi_E(x) dx,$$
where the last two integrals are the upper and lower Darboux integrals respectively.
By the Darboux criterion for existence of the integral (Theorem 3) the measure
$\mu(E)$ of a set is defined if and only if these lower and upper integrals are equal.
By the theorem of Darboux (Theorem 2 of Sect. 11.1) they are the limits of the
upper and lower Darboux sums of the function $\chi_E$ corresponding to partitions $P$
of $I$. But by definition of $\chi_E$ the lower Darboux sum is the sum of the volumes
of the intervals of the partition $P$ that are entirely contained in $E$ (the volume of a
polyhedron inscribed in $E$), while the upper sum is the sum of the volumes of the
intervals of $P$ that intersect $E$ (the volume of a circumscribed polyhedron). Hence
$\mu(E)$ is the common limit as $\lambda(P) \to 0$ of the volumes of polyhedra inscribed in
and circumscribed about $E$, in agreement with the accepted idea of the volume of
simple solids $E \subset \mathbb{R}^n$.
For $n = 1$ content is usually called length, and for $n = 2$ it is called area.
\begin{remark}{3} Let us now explain why the measure $\mu(E)$ introduced in Definition 3 is
sometimes called Jordan measure.
\end{remark}
\begin{definition}{4} A set $E \subset \mathbb{R}^n$ is a set of measure zero in the sense of Jordan or a set
of content zero if for every $\varepsilon > 0$ it can be covered by a finite system of intervals
$I_1, \dots, I_k$ such that $\sum_{i=1}^k |I_i| < \varepsilon$.
\end{definition}
Compared with measure zero in the sense of Lebesgue, a requirement that the
covering be finite appears here, shrinking the class of sets of Lebesgue measure
zero. For example, the set of rational points is a set of measure zero in the sense of
Lebesgue, but not in the sense of Jordan.
In order for the least upper bound of the contents of polyhedra inscribed in a
bounded set $E$ to be the same as the greatest lower bound of the contents of poly-
hedra circumscribed about $E$ (and to serve as the measure $\mu(E)$ or content of $E$), it
is obviously necessary and sufficient that the boundary $\partial E$ of $E$ have measure $0$ in
the sense of Jordan. That is the motivation for the following definition.
\begin{definition}{5} A set $E$ is Jordan-measurable if it is bounded and its boundary has
Jordan measure zero.
\end{definition}
As Remark 2 shows, the class of Jordan-measurable subsets is precisely the class
of admissible sets introduced in Definition 1. That is the reason the measure $\mu(E)$